Partisan Bias, defined as $\text{Bias} = S(0.50) - 0.50$ where $S(v)$ is the Democratic seat share
at statewide vote share $v$, measures whether the two parties are treated symmetrically by a
redistricting plan at the electoral pivot. It is estimated via the uniform swing model: each
district's vote share is shifted by $\delta = 0.50 - \bar{v}_D$ to bring the statewide mean to 50\%,
and the resulting seat share is computed. The symmetry standard---the framework within which Partisan
Bias is most often applied---requires that the seats-votes relationship be symmetric: if party A would
win $S\%$ of seats with $V\%$ of votes, then party B should win $S\%$ of seats if it also wins $V\%$
of votes. Zero Partisan Bias is a necessary (but not sufficient) condition for this symmetry.

We find that \textsc{bisect} algorithmic maps produce $|\text{Bias}| < 0.02$ across NC ($k=14$),
WI ($k=8$), and TX ($k=38$)---less than one seat in the NC delegation---compared to enacted NC
Bias $\approx -0.07$ (Republican advantage of approximately one full seat). Compactness-driven
redistricting produces near-zero Partisan Bias as a byproduct, because the symmetric treatment
of geographic space implies approximately symmetric treatment of partisan communities.

The paper documents the uniform swing assumption and its limitations. For Texas, with a high
proportion of safe seats, uniform swing overestimates Bias magnitude by approximately 15\%;
a heteroskedastic-swing correction narrows the gap but does not change the qualitative
conclusion. Partisan Bias has been cited in NC, PA, and OH expert testimony as part of the
symmetry standard framework; no court has adopted a specific Bias threshold as a constitutional
standard.
